% !TEX spellcheck = en_US
% !TEX spellcheck = LaTeX
\documentclass[letterpaper,10pt,english]{article}
\usepackage{%
amsfonts,%
amsmath,%
amssymb,%
amsthm,%
babel,%
bbm,%
%biblatex,%
caption,%
centernot,%
color,%
enumerate,%
%enumitem,%
epsfig,%
epstopdf,%
etex,%
fancybox,%
framed,%
fullpage,%
%geometry,%
graphicx,%
hyperref,%
latexsym,%
mathptmx,%
mathtools,%
multicol,%
paralist,%
pgf,%
pgfplots,%
pgfplotstable,%
pgfpages,%
proof,%
psfrag,%
%subfigure,%
tcolorbox,%
tfrupee,%
tikz,%
times,%
%ulem,%
url,%
xcolor,%
mathpazo,%
}

\definecolor{shadecolor}{gray}{.95}%{rgb}{1,0,0}
\usepackage[margin=1in,top=0.75in]{geometry}
\usepackage[mathscr]{eucal}
\usepgflibrary{shapes}
\usepgfplotslibrary{fillbetween}
\usetikzlibrary{%
arrows,%
backgrounds,%
chains,%
decorations.pathmorphing,% /pgf/decoration/random steps | erste Graphik
decorations.text,% 
matrix,%
positioning,% wg. " of "
fit,%
patterns,%
petri,%
plotmarks,%
scopes,%
shadows,%
shapes.misc,% wg. rounded rectangle
shapes.arrows,%
shapes.callouts,%
shapes%
}

%\pgfplotsset{compat=newest} %<------ Here
\pgfplotsset{compat=1.11} %<------ Or use this one

\theoremstyle{plain}
\newtheorem{thm}{Theorem}[section]
\newtheorem{lem}[thm]{Lemma}
\newtheorem{prop}[thm]{Proposition}
\newtheorem{cor}[thm]{Corollary}
\newtheorem{clm}[thm]{Claim}

\theoremstyle{definition}
\newtheorem{defn}[thm]{Definition}
\newtheorem{conj}[thm]{Conjecture}
\newtheorem{exmp}[thm]{Example}
\newtheorem{axiom}[thm]{Axiom}
\newtheorem{assum}[thm]{Assumption}
\newtheorem{exerc}[thm]{Exercise}
\newtheorem*{defin}{Definition}

\theoremstyle{remark}
\newtheorem{rem}{Remark}
\newtheorem{note}{Note}

\newcommand{\iid}{\emph{i.i.d.}\ }
\newcommand{\Cov}{\operatorname{Cov}}
%\newcommand{\det}{\operatorname{det}}
%\newcommand{\Real}{\mathbb{R}}
\newcommand{\tr}{\operatorname{tr}}
\newcommand{\Var}{\operatorname{Var}}
\newcommand{\cov}{\operatorname{cov}}

%\renewcommand{\proof}[1]{\begin{proof}#1\end{proof}}
\newcommand{\EQ}[1]{\begin{equation*}#1\end{equation*}}
\newcommand{\EQN}[1]{\begin{equation}#1\end{equation}}
\newcommand{\eq}[1]{\begin{align*}#1\end{align*}}
\newcommand{\meq}[2]{\begin{xalignat*}{#1}#2\end{xalignat*}}
\newcommand{\norm}[1]{\left\lVert#1\right\rVert}
\newcommand{\inner}[1]{\left\langle #1\right\rangle}
\newcommand{\abs}[1]{\left\lvert#1\right\rvert}
\newcommand{\expect}[1]{\mathbb{E}\left[{#1}\right]}
\newcommand{\prob}[1]{\mathbb{P}\left[{#1}\right]}
\newcommand{\given}{\; \big\vert \;} 
\newcommand{\set}[1]{\left\{#1\right\}} 
\newcommand{\indicator}[1]{1_{\set{#1}}} 
\newcommand{\Ind}[1]{\mathbbm{1}_{#1}} 
\newcommand{\SetIn}[1]{\mathbbm{1}_{\set{#1}}} 
\newcommand{\red}[1]{\textcolor{red}{#1}} 
\newcommand{\floor}[1]{\left\lfloor#1\right\rfloor}
\newcommand{\ceil}[1]{\left\lceil#1\right\rceil}


\newcommand{\C}{\mathbb{C}}
\newcommand{\D}{\mathbb{D}}
\newcommand{\E}{\mathbb{E}}
\newcommand{\F}{\mathbb{F}}
\newcommand{\N}{\mathbb{N}}
\renewcommand{\P}{\mathbb{P}}
\newcommand{\Q}{\mathbb{Q}}
\newcommand{\R}{\mathbb{R}}
\newcommand{\Z}{\mathbb{Z}}

\newcommand{\bA}{\mathbf{A}}
\newcommand{\bI}{\mathbf{I}}
\newcommand{\bU}{\mathbf{1}}
\newcommand{\bx}{\mathbf{x}}
\newcommand{\bX}{\mathbf{X}}
\newcommand{\bZ}{\mathbf{Z}}


\newcommand{\cA}{\mathcal{A}}
\newcommand{\cB}{\mathcal{B}}
\newcommand{\cC}{\mathcal{C}}
\newcommand{\cF}{\mathcal{F}}
\newcommand{\cG}{\mathcal{G}}
\newcommand{\cM}{\mathcal{M}}
\newcommand{\cN}{\mathcal{N}}
\newcommand{\cP}{\mathcal{P}}
\newcommand{\cT}{\mathcal{T}}
\newcommand{\cX}{\mathcal{X}}
\newcommand{\cY}{\mathcal{Y}}

\newcommand{\sA}{\mathscr{A}}
\newcommand{\sB}{\mathscr{B}}
\newcommand{\sC}{\mathscr{C}}
\newcommand{\sE}{\mathscr{E}}
\newcommand{\sF}{\mathscr{F}}
\newcommand{\sG}{\mathscr{G}}
\newcommand{\sH}{\mathscr{H}}
\newcommand{\sL}{\mathscr{L}}
\newcommand{\sS}{\mathscr{S}}
\newcommand{\sT}{\mathscr{T}}
\newcommand{\sX}{\mathscr{X}}
\newcommand{\sY}{\mathscr{Y}}
\newcommand{\sZ}{\mathscr{Z}}

\newcommand{\tS}{\tilde{S}}
% Debug
\newcommand{\todo}[1]{\begin{color}{blue}{{\bf~[TODO:~#1]}}\end{color}}

% a few handy macros

\renewcommand{\le}{\leqslant}
\renewcommand{\ge}{\geqslant}
\newcommand\matlab{{\sc matlab}}
\newcommand{\goto}{\rightarrow}
\newcommand{\bigo}{{\mathcal O}}
%\newcommand{\half}{\frac{1}{2}}
%\newcommand\implies{\quad\Longrightarrow\quad}
\newcommand\reals{{{\rm l} \kern -.15em {\rm R} }}
\newcommand\complex{{\raisebox{.043ex}{\rule{0.07em}{1.56ex}} \hskip -.35em {\rm C}}}


% macros for matrices/vectors:

% matrix environment for vectors or matrices where elements are centered
\newenvironment{mat}{\left[\begin{array}{ccccccccccccccc}}{\end{array}\right]}
\newcommand\bcm{\begin{mat}}
\newcommand\ecm{\end{mat}}

% matrix environment for vectors or matrices where elements are right justifvied
\newenvironment{rmat}{\left[\begin{array}{rrrrrrrrrrrrr}}{\end{array}\right]}
\newcommand\brm{\begin{rmat}}
\newcommand\erm{\end{rmat}}

% for left brace and a set of choices
%\newenvironment{choices}{\left\{ \begin{array}{ll}}{\end{array}\right.}
\newcommand\when{&\text{if~}}
\newcommand\otherwise{&\text{otherwise}}
% sample usage:
%\delta_{ij} = \begin{choices} 1 \when i=j, \\ 0 \otherwise \end{choices}


% for labeling and referencing equations:
\newcommand{\eql}{\begin{equation}\label}
\newcommand{\eqn}[1]{(\ref{#1})}
% can then do
%\eql{eqnlabel}
%...
%\end{equation}
% and refer to it as equation \eqn{eqnlabel}.
% some useful macros for finite difference methods:
\newcommand\unp{U^{n+1}}
\newcommand\unm{U^{n-1}}
% for chemical reactions:
\newcommand{\react}[1]{\stackrel{K_{#1}}{\rightarrow}}
\newcommand{\reactb}[2]{\stackrel{K_{#1}}{~\stackrel{\rightleftharpoons}
 {\scriptstyle K_{#2}}}~}
\makeatletter
\def\th@plain{%
\thm@notefont{}% same as heading font
\itshape % body font
}
\def\th@definition{%
\thm@notefont{}% same as heading font
\normalfont % body font
}
\makeatother
\date{}


\title{Lecture-07: Random Processes}
\author{}

\begin{document}
\maketitle

%%%%%%%%%%%%%%%%%%%%%%%%%%%%%%%%%%%%%%%%%%%%%%%%%%%%%%%
\section{Introduction} 
%%%%%%%%%%%%%%%%%%%%%%%%%%%%%%%%%%%%%%%%%%%%%%%%%%%%%%%
\begin{rem}
For an arbitrary index set $T$, and a real-valued function $x \in \R^T$, 
the projection operator $\pi_t:\R^T\to\R$ maps $x\in\R^T$ to $\pi_t(x) = x_t$. 
\end{rem}
\begin{defn}[Random process]
Let $(\Omega, \sF, P)$ be a probability space. 
For an arbitrary index set $T$ and state space $\sX \subseteq \R$, 
a map $X : \Omega \to \sX^T$ is called a \textbf{random process} if the projections $X_t:\Omega\to\sX$ defined by $\omega \mapsto X_t(\omega) \triangleq (\pi_t\circ X)(\omega)$ are random variables on the given probability space. 
\end{defn}
\begin{defn} 
For each outcome $\omega \in \Omega$, we have a function $X(\omega): T \mapsto \sX$ called the \textbf{sample path} or the \textbf{sample function} of the process $X$. %, 
%also written as the collection  
%\EQ{
%X(\omega) \triangleq %(X(\omega, t) \in \sX: t \in T) = 
%(X_t(\omega) \in \sX: t \in T).
%} 
\end{defn}
\begin{rem} 
A random process $X$ defined on probability space $(\Omega, \sF, P)$ with index set $T$ and state space $\sX\subseteq\R$, can be thought of as 
\begin{enumerate}[(a)]
\item a map $X: \Omega \times T \to \sX$, 
\item a map $X: T \to \sX^\Omega$, i.e. a collection of random variables $X_t: \Omega \to \sX$ for each time $t \in T$, 
\item a map $X: \Omega \to \sX^T$, i.e. a collection of sample functions $X(\omega): T \to \sX$ for each random outcome $\omega \in \Omega$. 
\end{enumerate}
%We will adopt the view that a random process is collection of sample functions.
\end{rem}

\subsection{Classification}
State space $\sX$ can be countable or uncountable, corresponding to discrete or continuous valued process.  
If the index set $T \subseteq \R$ is countable, the stochastic process is called \textbf{discrete}-time stochastic process or random sequence. 
When the index set $T$ is uncountable, it is called \textbf{continuous}-time stochastic process. 
The index set $T$ doesn't have to be time, if the index set is space, and then the stochastic process is spatial process. 
When $T = \R^n \times [0, \infty)$, stochastic process $X$ is a spatio-temporal process. 
\begin{shaded}
\begin{exmp}
We list some examples of each such stochastic process. 
\begin{enumerate}[i\_]
\item Discrete random sequence: brand switching, discrete time queues, number of people at bank each day.
\item Continuous random sequence: stock prices, currency exchange rates, waiting time in queue of $n$th arrival, workload at arrivals in time sharing computer systems.
\item Discrete random process:  counting processes, population sampled at birth-death instants, number of people in queues.
\item Continuous random process: water level in a dam, waiting time till service in a queue, location of a mobile node in a network.
\end{enumerate}
\end{exmp}
\end{shaded}
%%%%%%%%%%%%%%%%%%%%%%%%%%%%%%%%%%%%%%%%%%%%%%%
\subsection{Measurability}
%%%%%%%%%%%%%%%%%%%%%%%%%%%%%%%%%%%%%%%%%%%%%%%
For random process $X:\Omega\to \sX^T$ defined on the probability space $(\Omega, \sF, P)$, 
the projections $X_t \triangleq \pi_t\circ X$ are $\sF$-measurable random variables. 
Therefore, the set of outcomes $A_{X_t}(x) \triangleq X_t^{-1}(-\infty, x] \in \sF$ for all $t \in T$ and $x \in \R$.  
\begin{defn} 
A random map $X:\Omega\to\sX^T$ is called $\sF$-\textbf{measurable} and hence a random process,  
if the set of outcomes $A_{X_t}(x) = X_t^{-1}(-\infty, x]\in \sF$ for all $t \in T$ and $x\in \R$. 
\end{defn} 
\begin{defn}
The \textbf{event space generated by a random process} $X: \Omega\to\sX^T$ defined on a probability space $(\Omega,\sF,P)$ is given by 
\EQ{
\sigma(X) \triangleq \sigma(A_{X_t}(x): t\in T, x \in \R). 
}
\end{defn}
\begin{defn}
For a random process $X:\Omega \to \sX^T$ defined on the probability space $(\Omega,\sF,P)$,  
we define the projection of $X$ onto components $S \subseteq T$ as the random vector $X_S: \Omega \to \sX^S$, where $X_S \triangleq (X_s: s \in S)$.
\end{defn}
\begin{rem}
Recall that $\pi_t^{-1}(-\infty, x] = \bigtimes_{s \in T}(-\infty, x_s]$ where $x_s = x$ for $s = t$ and $x_s = \infty$ for all $s\neq t$. 
The $\sF$-measurability of process $X$ implies that for any countable set $S \subseteq T$, 
we have $A_{X_S}(x_S) \triangleq \cap_{s \in S}A_{X_s}(x_s) \in \sF$ for $x_S \in \sX^S$. 
\end{rem}
\begin{rem}
We can define $A_X(x) \triangleq \cap_{t \in T}A_{X_t}(x_t)$ for any $x \in \R^T$. 
However, $A_X(x)$ is guaranteed to be an event only when $S \triangleq \set{t \in T: \pi_t(x) < \infty}$ is a countable set. 
In this case, 
\EQ{
A_X(x) = \cap_{t \in T}A_{X_t}(x_t) = \cap_{s \in S}A_{X_s}(x_s) = A_{X_S}(x_S) \in \sF.    
}
%That is, when $y = \pi_S^{-1}(x_S) \in \R^T$ such that $\pi_s(y) = x_s$ for $s \in S$ and $\pi_t(y) = \infty$ for $t \notin S$.
\end{rem}
%For any finite subset $S \subseteq T$ and $x_S\in\R^S$, 
%we define 
%\begin{xalignat*}{2}
%&A_t(x_t) \triangleq \set{\omega\in\Omega: \pi_t\circ X \le x_t} = (\pi_t\circ X)^{-1}(-\infty, x_t],&
%&A_S(x_S) \triangleq \cap_{s\in S}A_s(x_s). 
%\end{xalignat*}
%
%For any finite subset $S \subseteq T$ and real vector $x \in \R^T$ such that $x_t = \infty$ for any $t \notin S$, 
%we define a set $A(x) = \set{y \in \R^T: y_t \le x_t} = \bigtimes_{t \in T}(-\infty, x_t]$. 
%Then, the measurability of the random process $X$ implies that for any such set $A(x)$, 
%we have 
%\EQ{
%\set{\omega \in \Omega: X_t(\omega) \le x_t, t \in T} 
%= \cap_{t \in T}X_s^{-1}(-\infty, x_t] 
%%= \cap_{s \in S}X_s^{-1}(-\infty, x_s]\cap_{t \notin S}X_t^{-1}(\R) 
%= X^{-1}\bigtimes_{t \in T}(-\infty, x_t] = X^{-1}(A(x)) \in \sF.
%}
%%where we have assumed $x_t = \infty$ for $t \in T\setminus S$. 
%
%\begin{rem}
%Realization of random process at each $t \in T$, is a random variable defined on the probability space $(\Omega, \sF, P)$ such that  $X_t: \Omega \to \sX$.  
%%\EQ{
%%X_t(\omega) \triangleq X(\omega, t) \in \sX, \text{ and }X_t = (X_t(\omega) \in \sX: \omega \in \Omega).
%%} 
%This follows from the fact that  for any $t \in T$ and $x_t \in \R$, we can take Boreal measurable sets $\bigtimes (-\infty, x_t]\bigtimes_{s \neq t}\R$. 
%Then, $X^{-1}(A(y)) = X_{t}^{-1}(-\infty, x_t] \in \sF$. 
%\end{rem}
%\begin{rem}
%The random process $X$ can be thought of as a collection of random variables $X = (X_t \in \sX^\Omega: t \in T)$ or an ensemble of sample paths $X = (X(\omega) \in \sX^T: \omega \in \Omega)$. 
%Recall that $\sX^T$ is set of all functions from the index set $T$ to state space $\sX$. 
%\end{rem}
%%each defined on the same probability space $(\Omega, \sF, P)$ is called a \textbf{random process} 
%%\end{defn}
\begin{shaded}
\begin{exmp}[Bernoulli sequence] 
\label{exmp:BernoulliSeqDefn}
%Consider a sample space of collection of infinite bi-variate sequences of successes (S) and failures (F) defined by $\Omega = \set{S, F}^\N$ and the index set $\N$.  
%An outcome $\omega \in \Omega$ is an infinite sequence $\omega = (\omega_1, \omega_2, \dots)$ such that $\omega_n  = \pi_n(\omega) \in \set{S, F}$ for each $n \in \N$. 
Consider a sample space $\set{H,T}^\N$. 
We define a mapping $X: \Omega \to \set{0,1}^\N$ such that $X_n(\omega) = \SetIn{H}(\omega_n) = \SetIn{\omega_n = H}$. %$X(\omega) = (\SetIn{S}(\omega_1), \SetIn{S}(\omega_2), \dots)$. 
%That is, we have 
%\begin{xalignat*}{2}
%& X_n(\omega) = \SetIn{S}(\omega_n), %&&X_n =(\indicator{\omega_n = S}: \omega_n \in \set{S, F}),&
%&X(\omega) = (\SetIn{S}(\omega_n): n \in \N).
%\end{xalignat*} 
The map $X$ is an $\sF$-measurable random sequence, if each $X_n: \Omega \to \set{0,1}$ is a bi-variate $\sF$-measurable random variable on the probability space $(\Omega,\sF,P)$. 
Therefore, the event space $\sF$ must contain the event space generated by sequence of events $E\in\sF^\N$ defined by $E_n \triangleq \set{\omega \in \Omega: X_n(\omega)=1} = \set{\omega \in \Omega: \omega_n = H} \in \sF$ for all $n\in \N$. 
That is, 
\EQ{
\sigma(X) = \sigma(E) = \sigma(\set{E_n: n \in \N}). 
}  
%Hence, we can write the map $X$ as the collection of random variables $X: \N \to \set{0,1}^\Omega$ or the collection of sample paths $X: \Omega \to \set{0,1}^\N$.
\end{exmp}
\end{shaded}
\subsection{Distribution}
%\subsection{Finite dimensional distribution} 
\begin{defn} 
For a random process $X: \Omega \to \sX^T$ defined on the probability space $(\Omega,\sF, P)$, 
we define a \textbf{finite dimensional distribution} $F_{X_S}: \R^S \to [0,1]$ for a finite $S \subseteq T$ 
by 
\EQ{
F_{X_S}(x_S) \triangleq P(A_{X_S}(x_S)),\quad x_S \in \R^S. 
}
\end{defn}
\begin{shaded}
\begin{exmp}
Consider a probability space $(\Omega,\sF, P)$ defined by the sample space $\Omega = \set{H,T}^\N$, the event space $\sF \triangleq \sigma(E)$ where $E_n = \set{\omega \in \Omega: \omega_n = H}$ for $n\in\N$, 
and the probability measure $P: \sF \to [0,1]$ defined by 
\EQ{
P(\cap_{i \in F}E_i) = p^{\abs{F}},\text{ for all finite }F \subseteq \N.
}
Let $X: \Omega\to\set{0,1}^\N$ defined as $X_n(\omega) = \Ind{E_n}(\omega)$ for all outcomes $\omega \in \Omega$ and $n \in \N$. 
For this random sequence, we can obtain the finite dimensional distribution $F_{X_S}: \R^S \to [0,1]$ for any finite $S \subseteq T$ and $x \in \R^S$ in terms of $I_0(x) \triangleq \set{i \in S: x_i < 0}$ and $I_1(x) \triangleq \set{i \in S: x_i \in [0,1)}$, 
as 
\EQN{
\label{eqn:FDDBernoulli}
F_{X_S}(x) = 
\begin{cases}
1,& I_0(x) \cup I_1(x) = \emptyset,\\
(1-p)^{\abs{I_1(x)}}, & I_0(x) = \emptyset, I_1(x) \neq \emptyset,\\
0, & I_0(x) \neq \emptyset.
\end{cases}
}
\end{exmp}
\end{shaded}
To define a measure on a random process, we can either put a measure on subsets of sample paths $(X(\omega) \in \R^T: \omega \in \Omega)$, 
or equip the collection of random variables $(X_t \in \R^\Omega: t \in T)$ with a joint measure. 
%We adopt the second approach, and 
Either way, we are interested in identifying the joint distribution $F: \R^T \to [0,1]$. 
To this end, for any $x \in \R^T$, we need to know
\EQ{
F_X(x) \triangleq P\left(\displaystyle {\bigcap_{t \in T}\{\omega \in \Omega: X_t(\omega) \le x_t\}}\right) 
= P(\bigcap_{t \in T}X_t^{-1}(-\infty, x_t]) 
= P(A_X(x)).
%= P \circ X^{-1}\bigtimes_{t \in T}(-\infty, x_t].
}
First of all, we don't know whether $A_X(x)$ is an event when $T$ is uncountable. 
Though, we can verify that $A_X(x) \in \sF$ for $x \in \R^T$ such that $\set{t \in T: x_t < \infty}$ is countable. 
Second, even for a simple independent process with countably infinite $T$, 
%When the index set $T$ is infinite, 
any function of the above form would be zero if $x_t$ is finite for all $t \in T$. 
Therefore, we only look at the values of $F_X(x)$ for $x \in \R^T$ where $\set{t \in T: x_t < \infty}$ is finite. 
% indices $t$ in a finite set $S$ and $x_t = \infty$  for all $t \notin S$. 
That is, for any finite set $S \subseteq T$, we focus on the events $A_S(x_S)$ and their probabilities. 
However, these are precisely the finite dimensional distributions.  
%product sets of the form 
%\EQ{
%A_S(x_S) \triangleq \bigtimes_{s \in S}(-\infty, x_s]\bigtimes_{s \notin S}\R.  
%}
%That is, $A(x) = A_S(x_S)$ where $x = \pi_S^{-1}(x_S) \in \R^T$ such that $x_t = \infty$ for $t \notin S$. 
%Recall that by definition of measurability, $X^{-1}(A(x)) \in \sF$, and hence $P\circ X^{-1} A(x)$ is well defined. 
%This leads to finite-dimensional distributions as defined below. 
%\begin{defn}[Finite dimensional distribution] 
%We can define a \textbf{finite dimensional distribution} for any finite set $S \subseteq T$ and $x_S \in \R^S$, 
%\EQ{
%F_{X_S}(x_S) \triangleq P\left(\displaystyle {\bigcap_{s \in S}\{\omega \in \Omega: X_s(\omega) \le x_s\}}\right) = P(\bigcap_{s \in S}X_s^{-1}(-\infty, x_s]).
%}
%\end{defn}
Set of all finite dimensional distributions of the stochastic process $X: \Omega \to \sX^T$ characterizes its distribution completely. 
\begin{shaded}
\begin{exmp}
Consider a probability space $(\Omega,\sF, P)$ defined by the sample space $\Omega = \set{H,T}^\N$ and the event space $\sF \triangleq \sigma(E)$ where $E_n = \set{\omega \in \Omega: \omega_n = H}$ for all $n\in\N$. 
Let $X: \Omega\to\set{0,1}^\N$ defined as $X_n(\omega) = \Ind{E_n}(\omega)$ for all outcomes $\omega \in \Omega$ and $n \in \N$. 
For this random sequence, if we are given the finite dimensional distribution $F_{X_S}: \R^S \to [0,1]$ for any finite $S \subseteq T$ and $x \in \R^S$ in terms of sets $ I_0(x) \triangleq \set{i \in S: x_i < 0}$ and $I_1(x) \triangleq \set{i \in S: x_i \in [0,1)}$, 
as defined in Eq.~\eqref{eqn:FDDBernoulli}. 
Then, we can find the probability measure $P: \sF \to [0,1]$ is given by 
\EQ{
P(\cap_{i \in F}E_i) = p^{\abs{F}},\text{ for all finite }F \subseteq \N.
}

\end{exmp}
\end{shaded}
\subsection{Independence}
%Recall the following definitions given the probability space $(\Omega, \sF, P)$. 
%Events $(A_i \in \sF: i \in [n]) \in \sF$ are independent if 
%$
%P(\cap_{i=1}^nA_i) = \prod_{i=1}^nP(A_i).
%$
%A collection of events $(A_t\in\sF: t \in T)$ is independent if $(A_s: s\in S)$ are independent for all finite $S\subseteq T$.   
%\begin{defn}
%Event spaces $(\sG_i \subseteq \sF: i \in [n])$ are independent if events $G_i \in \sG_i$ for all $i \in [n]$, are independent. 
%That is,
%\EQ{
%P(\cap_{i=1}^nG_i) = \prod_{i=1}^nP(G_i), \text{ for all } G_i \in \sG_i, i \in [n].
%} 
%\end{defn}
%\begin{defn}
%Random vector $X:\Omega\to\R^n$ is independent if event spaces $\sigma(X_1), \dots, \sigma(X_n)$ are independent. 
%That is, for all $x \in \R^n$, we have 
%\EQ{
%F_X(x) = P(\cap_{i=1}^nX_i^{-1}(-\infty, x_i]) = \prod_{i=1}^n P\circ X_i^{-1}(-\infty, x_i] = \prod_{i=1}^nF_{X_i}(x_i).
%}
%That is, independence of random vector is equivalent to factorization of joint distribution function into product of individual marginal distribution functions.
%\end{defn}
%\begin{defn}
%A collection of event spaces $(\sG_t \subseteq \sF: t \in T)$ is independent if any finite subcollection $(\sG_t: t \in S)$ for finite $S\subseteq T$ is independent. 
%\end{defn}
\begin{defn} 
A random process is \textbf{independent} if the collection of event spaces $(\sigma(X_t): t \in T)$ is independent. 
That is, for all $x_S \in \R^S$, we have 
\EQ{
F_{X_S}(x_S) = P(\cap_{s \in S}\set{X_s \le  x_s}) 
= \prod_{s \in S}P\set{X_s \le x_s} 
= \prod_{s \in S}F_{X_s}(x_s). 
}
That is, independence of a random process is equivalent to factorization of any finite dimensional distribution function into product of individual marginal distribution functions.
\end{defn}
%A stochastic process $X$ is said to be \textbf{independent} if for all finite subsets $S \subseteq T$, the finite collection of events $\set{\set{X_s \le x_s}: s \in S}$ are independent.  
\begin{shaded}
\begin{exmp}
Consider a probability space $(\Omega,\sF, P)$ defined by the sample space $\Omega = \set{H,T}^\N$, the event space $\sF \triangleq \sigma(E)$ where $E_n = \set{\omega \in \Omega: \omega_n = H}$ for all $n\in\N$, 
and the probability measure $P: \sF \to [0,1]$ defined by 
\EQ{
P(\cap_{i \in F}E_i) = p^{\abs{F}},\text{ for all finite }F \subseteq \N.
}
Then, we observe that the random sequence $X: \Omega\to\set{0,1}^\N$ defined by $X_n(\omega) \triangleq \Ind{E_n}(\omega)$ for all outcomes $\omega \in \Omega$ and $n \in \N$, is independent.
\end{exmp}
\end{shaded}

\begin{defn}
Two stochastic processes $X: \Omega \to \sX^{T_1},Y:\Omega\to \sY^{T_2}$ are \textbf{independent}, if the corresponding event spaces $\sigma(X), \sigma(Y)$ are independent. 
That is, for any $x \in \R^{S_1}, y \in \R^{S_2}$ for finite $S_1 \subseteq T_1, S_2\subseteq T_2$, 
the events $A_{S_1}(x) \triangleq \cap_{s\in S_1}X_s^{-1}(-\infty, x_s]$ and $B_{S_2}(y) \triangleq \cap_{s \in S_2}Y_s^{-1}(-\infty, y_s]$ are independent. 
That is, the joint finite dimensional distribution of $X$ and $Y$ factorizes, and 
\EQ{
P(A_{S_1}(x)\cap B_{S_2}(y)) 
= P(A_{S_1}(x))P(B_{S_2}(y)) 
= F_{X_{S_1}}(x)F_{Y_{S_2}}(y),\quad x \in \R^{S_1}, y \in \R^{S_2}.
}
\end{defn}
%Two stochastic process $X, Y$ for the common index set $T$ are \textbf{independent random processes} if for all %$m,n \in \N$ and 
%finite subsets $I, J \subseteq T$, 
%the following events $\set{X_i \le x_i, i \in I}$ and $\set{Y_j \le y_j, j \in J}$ are independent. %such that $|I| =m, |J| = n$, 
%That is, 
%\EQ{
%F_{X_I, X_J}(x_I, x_J)
%\triangleq P\left(\set{X_i \le x_i, i \in I}\cap\set{Y_j \le y_j, j \in J}\right) 
%= P\left(\cap_{i \in I}\set{X_i \le x_i}\right)P\left(\cap_{j \in J}\set{Y_j \le y_j}\right)
%= F_{X_I}(x_I)F_{X_J}(x_J).
%}
 
\end{document}

\subsection{Joint moments}
Simpler characterizations of a stochastic process $X$ are in terms of its moments. 
That is, the first moment such as mean, and the second moment such as correlations and covariance functions. 
\meq{3}{
&m_X(t) \triangleq \E X_t,&
&R_X(t,s) \triangleq \E X_tX_s,&
&C_X(t,s) \triangleq \E (X_t - m_X(t))(X_s-m_X(s)).
}
%\begin{shaded}
%\begin{exmp}
%Some examples of simple stochastic processes. 
%\begin{enumerate}[i\_]
%\item $X_t = A \cos 2\pi t$, where $A$ is random. 
%The finite dimensional distribution is given by 
%\EQ{
%F_{X_S}(x) = P\left(\left\{A\cos 2\pi s \leq x_s, s \in S\right\}\right). %P\left(\left\{A \leq \min_{s \in S\setminus\{(2k+1)\frac{\pi}{2}, k \in \Z\}}\frac{x_s}{\cos 2\pi_s}\right\}\right).
%}
%The moments are given by 
%\meq{3}{
%&m_X(t) = (\E A)\cos 2\pi t, && R_X(t,s) =  (\E A^2) \cos 2\pi t\cos 2\pi s,&& C_X(t,s) =\text{Var}(A) \cos 2\pi t\cos 2\pi s.
%}
%\item $X_t = \cos(2\pi t+ \Theta)$, where $\Theta$ is random and uniformly distributed between $(-\pi, \pi]$. 
%The finite dimensional distribution is given by 
%\EQ{
%F_{X_S}(x) = P\left(\left\{\cos(2\pi s + \Theta) \leq x_s, s \in S\right\}\right). %P\left(\left\{A \leq \min_{s \in S\setminus\{(2k+1)\frac{\pi}{2}, k \in \Z\}}\frac{x_s}{\cos 2\pi_s}\right\}\right).
%}
%The moments are given by 
%\meq{3}{
%&m_X = 0, && R_X(t,s) =  \frac{1}{2}\cos2\pi (t-s),&& C_X(t,s) = R_X(t,s).
%}
%\item $X_n = U^n$ for $n \in \N$, where $U$ is uniformly distributed in the open interval $(0,1)$.
%\item $Z_t = At +B$ where $A$ and $B$ are independent random variables. 
%\end{enumerate}
%\end{exmp}
%\end{shaded}
%
\begin{shaded}
\begin{exmp}[Bernoulli sequence] 
\label{exmp:BernoulliSeqIndep}
For the Bernoulli sequence $X$ defined in Example~\ref{exmp:BernoulliSeqDefn} on probability space $(\Omega,\sF,P)$, 
we take the probability of generating events be 
\EQ{
P(\cap_{i\in S}E_i) = p^{\abs{S}},\quad S \subseteq \N.
} 
This probability of events implies that finite dimensional distributions factorize and they are identical. 
In particular, we have 
\EQ{
P(\cap_{s\in S}\set{X_s = x_s}) = p^{\sum_{s \in S}x_s}(1-p)^{\abs{S}-\sum_{s\in S}x_s},\quad S \subseteq \N.  
}
Hence, the random sequence $X$ is independent and identically distributed with $P\set{X_n = 1} = p \in (0,1)$. 
For any sequence $x \in \set{0,1}^\N$, we have $P\set{X = x} = 0$. 
Let $q \triangleq (1-p)$, then the probability of observing $m$ heads and $r$ tails is given by $p^mq^r$. 
%\end{exmp}
%\end{shaded}
%\begin{shaded}
%\begin{exmp}[Bernoulli sequence] 
We can easily compute the mean, the auto-correlation, and the auto-covariance functions for the independent Bernoulli process defined in Example~\ref{exmp:BernoulliSeqIndep} as 
\meq{3}{
&m_X(n) = \E X_n = p, &&R_X(m,n)=\E X_mX_n = \E X_m\E X_n = p^2,&&C_x(m,n) = 0.
}
\end{exmp}
\end{shaded}
\end{document}

\section{Examples of Tractable Stochastic Processes}
In general, it is very difficult to characterize a stochastic process completely in terms of its finite dimensional distribution. 
However, we have listed few analytically tractable examples below, where we can completely characterize the stochastic process. 
\subsection{Independent and identically distributed processes}
Let $(X_t: t \in T)$ be an independent and identically distributed (\emph{i.i.d.}) random process, with a common distribution $F(x)$. 
Then, the finite dimensional distribution for this process for any finite $S \subseteq T$ can be written as 
\EQ{
F_S(x_S) = P\left(\bigcap_{s \in S}\set{X_s(\omega) \le x_s, s \in S}\right) = \prod_{s \in S}F(x_s).
}
It's easy to verify that the first and the second moments are independent of time indices. 
Since $X_t = X_0$ in distribution, 
\meq{3}{
&m_X = \E X_0 , && R_X = \E X_0^2, && C_X = \Var(X_0).
}

%\subsubsection{Independent increments process}
\subsection{Stationary processes}
A stochastic process $X$ is \textbf{stationary} if all finite dimensional distributions are shift invariant. 
That is, for any finite $n \in \N$ and $t > 0$, the random vectors $(X_{s_1}, \dots, X_{s_n})$ and $(X_{s_1+t}, \dots, X_{s_1+t})$ have the identical joint distribution for all $s_1 \le \dots \le s_n$.  
That is, for finite $S \subseteq T$ and $t > 0$, we have 
\EQ{
F_S(x_S) = P(\{X_s(\omega) \le x_s, s \in S\}) = P(\{X_{s+t}(\omega) \le x_s, s \in S\}) = F_{t+S}(x_S).
}
For a stationary stochastic process, all the existing moments are shift invariant. 
For Gaussian random processes, first and the second moment suffice to get any finite dimensional distribution. 
A \textbf{second order} stochastic process $X$ has finite auto-correlation $R_X(t,t) < \infty$ for all indices $t \in T$. 
This implies $R_X(t_1, t_2) < \infty$ by Cauchy-Schwartz inequality, 
and hence the mean, auto-correlation, and the auto-covariance functions are well defined and finite. 
Since $X_t = X_0$ and $(X_t,X_s) = (X_{t-s}, X_{0})$ in distribution, we have for a second order process $X$ 
\meq{3}{
&m_X = \E X_0 , && R_X(t, s) = R_X(t-s,0) = \E X_{t-s}X_0, && C_X(t-s, 0) = R_X(t-s, 0) - m_X^2.
}


A random process $X$ is \textbf{wide sense stationary} if 
\begin{enumerate}
\item $m_X(t) = m_X(t+s)$ for all $s,t \in T$, and 
\item $R_X(t,s) = R_x(t+u, s+u)$ for all $s,t,u \in T$. 
\end{enumerate} 
It follows that a second order stationary stochastic process $X$,  is wide sense stationary. 
A second order wide sense stationary process is not necessarily stationary. 
We can similarly define join stationarity and joint wide sense stationarity for two stochastic processes $X$ and $Y$.

\subsection{Markov processes}
A stochastic process $X$ is \textbf{Markov} if conditioned on the present state, future is independent of the past. 
We denote the history of the process until time $t$ as $\sF_t = \sigma(X_s, s \le t)$. 
That is, for any ordered index set $T$ containing any two indices $u > t$, we have  %such that $t < u$, % \in T$, 
\EQ{
P(\{X_u \le x_u\}\given \sF_t) = P(\{X_{u} \le x_u\}\given \sigma(X_t)).
}
The range of the process is called the \textbf{state space}. 
%For a stationary Markov process, we would have 
We next re-write the Markov property more explicitly for the process $X$. 
For all $x, y \in \sX$, finite set $S \subseteq T$ such that $\max S < t < u$, and $H_{s}  = \set{X_s \le x_s: s \in S} \in \sF_t$, we have 
\EQ{
P(\{X_u \le y\}\given H_S \cap \set{X_t \le x}) = P(\{X_{u} \le y\}\given \set{X_t \le x}).
}
When the state space $\sX$ is countable, we can write $H_S = \cap_{s \in S}\set{X_s = x_s}$ and the Markov property can be written as 
\EQ{
P(\{X_u = y\}\given H_S \cap \set{X_t = x}) = P(\{X_{u} = x_u\}\given \set{X_t = x}).
}
We will study this process in detail in coming lectures. 


\begin{shaded}
\begin{exmp}[Random Walk] 
Let $X = (X_n \in \R^d: n \in \N)$ be an be an independent (not necessarily identical) Bernoulli sequence. 
Let $S_0 = 0$ and $S_n \triangleq \sum_{i=1}^nX_i$, then the process $S = (S_n \in \R^d: n \in \N_0)$ is called a \textbf{random walk}. 
We can think of $S_n$ as the random location of a particle after $n$ steps, 
where the particle starts from origin and takes steps of size $X_i$ at the $i$th step. 

%From previous section, we know following properties of random walks. 
\begin{thm} 
For a random walk $(S_n: n \in \N)$ with independent step-size sequence $X$, the following are true. 
\begin{enumerate}[i\_]
\item The first two moments are $\E S_n = \sum_{i=1}^n\E X_i$ and $\Var[S_n] = \sum_{i=1}^n\Var[X_i]$. 
\item Random walk is non-stationary with independent increments. 
\item Random walk is a Markov sequence. 
\end{enumerate}
\end{thm}
\proof{
Results follow from the independence of the step-size sequence $X$. 
\begin{enumerate}[i\_]
\item Follows from the linearity of expectation and independence of step sizes. 
\item Since the mean is time dependent, random walk is non-stationary process. 
Independence of increments follows from the independence of step sizes. 
\item Given the historical event $H_{n-1} \triangleq \cap_{k=1}^{n-1}\set{S_k \le s_k}$ and the current state $\set{S_n \le s_n}$, 
we can write the conditional probability 
\eq{
P(\set{S_{n+1} \le s_{n+1}}\given H_{n-1}\cap\set{S_n \le s_n}) &= P(\set{X_{n+1} \le s_{n+1}-S_n}\given H_{n-1}\cap\set{S_n \le s_n})\\
&= P(\set{S_{n+1} \le s_{n+1}}\given \set{S_n \le s_n}). % = P\set{X_{n+1} = s_{n+1}-s_n} .
}
The equality in the second line follows from the independence of the step-size sequence. 
In particular, from the independence of $X_{n+1}$ from the collection $\sigma(S_0, X_1, \dots, X_n) = \sigma(S_0, S_1, \dots, S_n)$. 
For the countable state space $\sX$, an given the historical event $H_{n-1} \triangleq \cap_{k=1}^{n-1}\set{S_k = s_k}$ and the current state $\set{S_n = s_n}$, 
we can write the conditional probability 
\eq{
P(\set{S_{n+1} = s_{n+1}}\given H_{n-1}\cap\set{S_n = s_n}) &= P(\set{X_{n+1} = s_{n+1}-S_n}\given H_{n-1}\cap\set{S_n = s_n})\\
&= P(\set{S_{n+1} = s_{n+1}}\given \set{S_n = s_n}) = P\set{X_{n+1} = s_{n+1}-s_n} .
}
\end{enumerate}
}
\end{exmp}
\end{shaded}

\subsection{L\'evy processes}
A right continuous with left limits stochastic process $X = (X_t \in \R: t \in T \subseteq \R_+)$ with $X_0 = 0$ almost surely, is a \textbf{L\'evy process} if the following conditions hold. 
\begin{enumerate}[(L1)]
%\item $X_{0}=0$ almost surely.
\item The increments are independent. For any $ 0\leq t_{1}<t_{2}<\cdots <t_{n}<\infty$, $X_{t_{2}}-X_{t_{1}},X_{t_{3}}-X_{t_{2}},\ldots ,X_{t_{n}}-X_{t_{n-1}}$ are independent.
\item The increments are stationary. For any $s<t$,  $X_{t}-X_{s}$, is equal in distribution to $X_{t-s}$.
\item Continuous in probability. For any $\epsilon > 0$ and $t\geq 0$ it holds that $\lim _{h\rightarrow 0}P(|X_{t+h}-X_{t}|>\epsilon )=0$. 
\end{enumerate}

\begin{shaded}
\begin{exmp}
Two examples of L\'evy processes are Poisson process and Wiener process. 
The distribution of Poisson process at time $t$ is Poisson with rate $\lambda t$ and the distribution of Wiener process at time $t$ is zero mean Gaussian with variance $t$. 
\end{exmp}
\end{shaded}

\section{Conditional Expectation} 
Let $(\Omega, \sF, P)$ be the probability space.  
Let $X$ be a measurable random variable on this probability space denoted as $X \in \sF$, 
if the event $X^{-1}(-\infty, x] = \{\omega \in \Omega: X(\omega) \leq x\} \in \sF$ for each $x \in \R$.  
Let $\sE \subseteq \sF$ be a $\sigma$-algebra, then the \textbf{conditional expectation} of $X$ given $\sE$ is denoted $\E[X | \sE]$ and is a random variable $Y = \E[X | \sE]$ where
\begin{enumerate}[i\_] 
\item $Y \in \sE$,
\item for each event $A \in \sE$, we have $\E [X1_A] = \E[Y1_A]$. 
%\eq{
%\int_A X dP = \int_A Y dP. 
%}
\end{enumerate}
Intuitively, we think of the $\sigma$-algebra $\sE$ as describing the information we have. 
For each $A \in \sE$, we know whether or not $A$ has occurred. 
The conditional expectation $\E[X|\sE]$ is then the ``best guess'' of the value of $X$ given the information $\sE$. 
Let $X,Y$ be two random variables defined on this probability space. 
Then, the conditional expectation of $X$ given $Y$ is defined as 
\eq{
\E[X|Y] = \E[X | \sigma(Y)]. 
}
A random variable $X$ is \textbf{independent} of the $\sigma$-algebra $\sE$, if for all $x \in \R$ and $A \in \sE$, 
\eq{
\E[1_{\{X \leq x\}}1_A]= P\{X \leq x\}\cap A = P\{X \leq x\}P(A) = \E1_{\{X \leq x\}}\E1_A. 
}
\begin{lem}
Let $(\Omega, \sF, P)$ be a probability space with $\sE \subseteq \sF$ a $\sigma$-algebra. 
If $X \in \sE$ is a random variable, then $\E[X| \sE] = X$. 
\end{lem}
\begin{proof} 
First condition is true by hypothesis, and the second condition holds for any $A \in \sE$. 
\end{proof}
\begin{lem} 
Let $(\Omega, \sF, P)$ be a probability space with $\sE \subseteq \sF$ a $\sigma$-algebra. 
If $X \in \sF$ be a random variable independent of $\sE$.  
Then, $\E[X| \sE]  = \E[X]$. 
\end{lem}
\begin{proof}
This follows since $\E X \in \sE$ and the random variables $X$ and $1_{A}$ are independent for any $A \in \sE$, 
which implies 
\eq{
\E[X1_A] = \E X \E 1_A = \E [(\E X) 1_A]. 
}
\end{proof}
One can partition the state space $\R$ into measurable sets $E_1, E_2, \dots$ for the random variable $X$ defined on the given probability space. 
Then $\Omega_i \triangleq X^{-1}(E_i)$ is a partition of the sample space $\Omega$. 
Let $Y$ be a random variable defined as the partition index for the random variable $X$. 
That is, 
\eq{
Y = \sum_{i \in \N}i \cdot 1_{\{X \in E_i\}}. 
} 
Let $\sE \triangleq \sigma(\Omega_1, \Omega_2, \dots)$, then one can check that $Y \in \sE$ or $\sigma(Y) = \sE$. 
Hence, $\E[X|Y] = \E[X|\sigma(Y)] = \E[X|\sE]$. 
Clearly, $\E[X|Y]$ would be a function of $Y$ and since $Y$ takes countably many values, we have $Z = \E[X|Y]$ taking countably many values, with $Z_i = Z1_{\{Y = i\}}$ being a constant on the corresponding partition $\Omega_i$  of the sample space. 
One can compute this conditional expectation using joint distribution directly as 
\eq{
\E[X|Y=i] = \int_{\R}xdF_{X|Y=i}(x) = \frac{1}{P(\Omega_i)}\int_{E_i}xdF(x) = \frac{\E X1_{\Omega_i}}{P(\Omega_i)}
}
\begin{lem}
Suppose $\{\Omega_i: i \in \N\}$ be a countable partition of the sample space $\Omega$, and $\sE = \sigma(\Omega_1, \Omega_2, \dots)$ is the $\sigma$-field generated by this partition. 
Then, 
\eq{
\E[X | \sE] = \frac{\E[X1_{\Omega_i}]}{P(\Omega_i)} \text{ on } \Omega_i.
}
\end{lem}
\begin{proof}
It is easy to see that the RHS is constant on each partition $\Omega_i$ and hence is measurable with respect to $\sE$. 
Further, for each $\Omega_i \in \sE$, we have
\eq{
\int_{\Omega_i}\frac{\E[X 1_{\Omega_i}]}{P(\Omega_i)}dP = \E[X 1_{\Omega_i}] = \int_{\Omega_i}XdP. 
}
\end{proof}
\begin{cor} $P(A|B)P(B) = P(A\cap B)$.
\end{cor}
\begin{proof}
Taking $X = 1_A$ and $\sE = \{\emptyset, \Omega, B, B^c\}$, from the previous Lemma we get 
\eq{
P(A|B) = \E[1_A| 1_B]= \E[1_A|\sE] = \frac{\E[1_A1_B]}{P(B)} = \frac{P(A\cap B)}{P(B)}.
} 
\end{proof}

\begin{thm}[Bayes' Formula] 
For a $\sigma$-algebra $\sE \subseteq \sF$, and for any events $G \in \sE$ and $A \in \sF$, we have 
\eq{
P(G|A) = \frac{\E[1_GP(A| \sE)]}{\E P(A|\sE)}.
}
\end{thm}
\begin{proof}
It is easy to check that numerator is $\E 1_G\E[1_A|\sE] = \E[1_{A \cap G}|\sE]$. 
It suffices to show that $\E\E[1_A|\sE] = \E 1_A$, which follows from definition. 
\end{proof}
\begin{cor} 
For the countable partition $(\Omega_1, \Omega_2, \dots)$, if the $\sigma$-algebra $\sE = \sigma(\Omega_1, \Omega_2, \dots)$, then for any events $G \in \sE$ and $A \in \sF$, we have 
\eq{
P(\Omega_i|A) = \frac{P(A|\Omega_i)P(\Omega_i)}{\sum_{j \in \N}P(A|\Omega_j)P(\Omega_j)}.
}
\end{cor}
\begin{proof}
Result follows from the fact that $P(A|\sE) \in \sE$ and hence is a constant on each partition $\Omega_j$. % and $\sum_{j \in \N}1_{\Omega_j} = 1$.
\end{proof}

\subsection{Filtration}
A net of $\sigma$-algebras $\sF_{\bullet} = \{\sF_t \subseteq \sF: t \in T\}$ is called a \textbf{filtration} when the index set $T$ is totally ordered and the net is non-decreasing, that is for all $s \leqslant t \in T$ implies $\sF_s \subseteq \sF_t$. 
Consider a random process $X$ indexed by the ordered set $T$ on the probability space $(\Omega, \sF, P)$. 
The process $X$ is called \textbf{adapted} to the filtration $\sF_{\bullet}$, if for each $t \in T$, we have the random variable $X_t \in \sF_t$. 
For a random process $X$ with an ordered index set $T$, we can define a natural filtration $\sF_{\bullet} = \{\sF_t \subseteq \sF: t \in T\}$ indexed by $T$, where  $\sF_t \triangleq \sigma(X_s, s \leqslant t)$ is the information about the process till index $t$ and 
the process $X$ is adapted to its natural filtration by definition. 

If $X = (X_t: t \in T)$ is an independent process with the associated natural filtration $\sF_{\bullet}$, then for any $t > s$ and events $A \in \sF_s$, 
$X_t$ is independent of the event $A$. 
This is just a fancy way of saying $X_t$ is independent of $(X_u, u \le s)$. 
Hence, for any random variable $Y \in \sF_s$, we have 
\eq{
\E[\E[X_tY|\sF_s]] = \E[\E[X_t] Y] = \E X_t \E Y.
}

\end{document}